\begin{pastexam}[label={px:ntu-106-esoe-ds-01}]{臺大工科海洋所 106 學年度 第 1 題}
（程式設計）

\textbf{(15 \%)} A singly linked list originally stores $n$ nodes of integers not in order. It is
desired to convert this singly linked list into a special doubly linked list with a particular
dummy node, which is not seen from outside. The converted list has three different sections. The
first section contains all nodes with negative integers starting from the head. The second
section is the dummy node, whose value indicates the number of zeros encountered during the
conversion. The final section contains all nodes with positive integers to the end. There is no
need to sort the integers in the first and final sections.
\begin{enumerate}[label=\alph*.]
\item (5 \%) Write down the definitions of the singly and doubly linked list classes and nodes.
\item (10 \%) Write a program that accepts a singly linked list and returns the specified doubly
linked list.
\end{enumerate}
\end{pastexam}

\begin{pastsoln}
原卷規定程式題用 C++ 或 Java，這裡用 C++。

\textbf{想法。} 雙向串列一開始只放虛擬節點（dummy）。 把單向串列從頭走一遍：負數接在虛擬
節點前面，正數接在串列最後面，$0$ 不建節點，只把虛擬節點的值加 $1$。 虛擬節點永遠在負數和
正數之間，三段自然就分好了，不需要排序。 「外面看不到」的做法是把它設成 \texttt{private}，
並讓走訪用的 \texttt{first}、\texttt{after} 自動跳過它。 雙向串列的插入是資結書第~5.2~節
虛擬碼~5.6 的寫法。

\textbf{(a) 類別定義。}
\begin{minted}[fontsize=\small]{cpp}
struct SNode {                 // 單向串列的節點
    int val;
    SNode* next;
    SNode(int v, SNode* n = nullptr) : val(v), next(n) {}
};
class SList {                  // 單向串列
public:
    SNode* head = nullptr;
    void pushFront(int v) { head = new SNode(v, head); }
};

struct DNode {                 // 雙向串列的節點
    int val;
    DNode *prev, *next;
    DNode(int v) : val(v), prev(nullptr), next(nullptr) {}
};
class DList {                  // 負數 | 虛擬節點 | 正數
public:
    DList() { dummy = head = tail = new DNode(0); }  // 一開始只有虛擬節點
    ~DList();
    DNode* first() const {     // 外面看到的第一個節點（跳過虛擬節點）
        return head == dummy ? dummy->next : head;
    }
    DNode* after(DNode* p) const {   // p 的下一個（跳過虛擬節點）
        DNode* q = p->next;
        return q == dummy ? q->next : q;
    }
    int zeros() const { return dummy->val; }
    void insertBefore(DNode* pos, DNode* x);   // 把 x 接在 pos 前面
    void insertAfter(DNode* pos, DNode* x);    // 把 x 接在 pos 後面
    friend DList* convert(SList& s);
private:
    DNode *head, *tail, *dummy;  // 串列裡一定有 dummy
};
\end{minted}
\texttt{head}、\texttt{tail} 是整條雙向串列（含虛擬節點）的頭尾，所以串列永遠不是空的，插入時
只有「接在最前面」「接在最後面」兩種邊界要處理。

\textbf{(b) 轉換。}
\begin{minted}[fontsize=\small]{cpp}
DList::~DList() {
    for (DNode* p = head; p != nullptr; ) { DNode* q = p->next; delete p; p = q; }
}
void DList::insertBefore(DNode* pos, DNode* x) {
    x->next = pos; x->prev = pos->prev;
    if (pos->prev) pos->prev->next = x; else head = x;
    pos->prev = x;
}
void DList::insertAfter(DNode* pos, DNode* x) {
    x->prev = pos; x->next = pos->next;
    if (pos->next) pos->next->prev = x; else tail = x;
    pos->next = x;
}

DList* convert(SList& s) {
    DList* d = new DList();
    SNode* p = s.head;
    while (p != nullptr) {
        SNode* nxt = p->next;
        if (p->val < 0)
            d->insertBefore(d->dummy, new DNode(p->val));
        else if (p->val > 0)
            d->insertAfter(d->tail, new DNode(p->val));
        else
            d->dummy->val++;               // 0：只計數
        delete p;                          // 單向串列的節點用完就還
        p = nxt;
    }
    s.head = nullptr;
    return d;
}
\end{minted}

\textbf{正確性。} 迴圈的不變量：處理完單向串列的前 $t$ 個節點後，雙向串列是「這 $t$ 個裡的
負數（依原順序）、虛擬節點、這 $t$ 個裡的正數（依原順序）」，虛擬節點的值等於這 $t$ 個裡
$0$ 的個數。 第 $t + 1$ 個是負數時接在虛擬節點前面，也就是負數段的最後；是正數時接在
\texttt{tail} 後面，也就是正數段的最後；是 $0$ 時只改計數，三種情形都保持不變量。 迴圈結束
時 $t = n$，就是題目要的串列。

例如單向串列 $3, -1, 0, 5, -7, 0, 0, 2, -4, 9$ 轉換後是
$-1 \leftrightarrow -7 \leftrightarrow -4 \leftrightarrow [3] \leftrightarrow 3 \leftrightarrow 5
\leftrightarrow 2 \leftrightarrow 9$，方括號是虛擬節點，值 $3$ 是零的個數；從外面用
\texttt{first}、\texttt{after} 走訪只看到 $-1, -7, -4, 3, 5, 2, 9$。

\textbf{複雜度。} 每個節點處理一次，每次插入只改常數個指標，時間 $\Theta(n)$；新串列有
$n - z + 1$ 個節點（$z$ 是零的個數），舊節點邊走邊刪除，額外空間只多一個虛擬節點。 單向
節點沒有 \texttt{prev} 欄位，型別不同，不能直接把舊節點拿來用，所以是建新節點、刪舊節點。
\end{pastsoln}
